\chapter{Relational Automated Machine Learning Agents (\tool)}

Chương này trình bày kiến trúc và cơ chế hoạt động của \tool --- một khung học máy tự động dựa trên kiến trúc đa tác tử, hướng tới giải quyết các tác vụ dự đoán trên cơ sở dữ liệu quan hệ từ đầu đến cuối. Các hệ thống AutoML truyền thống như Auto-sklearn~\cite{feurer2015autosklearn} hay Auto-WEKA~\cite{kotthoff2017autoweka} tối ưu hóa trên tập dữ liệu dạng bảng đơn lẻ đã được trích xuất đặc trưng sẵn. \tool đối mặt với bài toán khác: tự động hóa toàn bộ quá trình từ cơ sở dữ liệu quan hệ thô $(\mathcal{T}, \mathcal{L})$ (Mục~\ref{sec:relational_data}) đến mô hình mạng nơ-ron đồ thị đã huấn luyện trên đồ thị thực thể quan hệ (Mục~\ref{sec:relation-graph}), theo định nghĩa của Relational Deep Learning~\cite{robinson2024relbench}.

So với các hệ thống AutoML truyền thống và các khung đa tác tử đa mục đích như MetaGPT~\cite{hong2023metagpt} hay AutoGen~\cite{wu2023autogen}, \tool khác biệt ở ba khía cạnh:
\begin{enumerate}
    \item \textbf{Điều phối dựa trên đồ thị trạng thái}: Pipeline đa giai đoạn được điều phối qua đồ thị trạng thái có phân nhánh điều kiện và tự phục hồi khi gặp lỗi, thay vì giao tiếp hội thoại tuyến tính giữa các tác tử.
    \item \textbf{Tri thức chuyên gia nhúng}: Mỗi tác tử nhận tri thức chuyên biệt về Relational Deep Learning qua system prompt --- bao gồm các mẫu thiết kế SQL cho bảng huấn luyện $T_{train}$ (Mục~\ref{sec:task_to_training_table}), phân loại bảng dữ kiện/chiều (Mục~\ref{sec:relational_data}), và xử lý đồ thị nhất quán theo thời gian (Mục~\ref{sec:time-consistent}).
    \item \textbf{Tối ưu siêu tham số không cần huấn luyện}: Hệ thống truy cập các nguồn tri thức ngoài (bài báo khoa học, bảng xếp hạng, cấu hình đã kiểm chứng) qua giao thức Model Context Protocol (MCP) để lựa chọn siêu tham số mà không cần thí nghiệm huấn luyện.
\end{enumerate}

Hệ thống phân rã quy trình thành sáu giai đoạn, mỗi giai đoạn do một tác tử chuyên biệt đảm nhiệm. Các tác tử chia sẻ một bộ nhớ trạng thái toàn cục (global pipeline state) để truyền tải thông tin giữa các pha mà không cần giao tiếp điểm-đến-điểm.

\section{Kiến Trúc Tổng Quát}
\label{sec:architecture}

Kiến trúc của \tool được tổ chức theo mô hình \textit{đồ thị trạng thái} (state graph), triển khai trên LangGraph --- một thư viện điều phối tác tử hỗ trợ đồ thị có hướng với các cạnh điều kiện (conditional edges), checkpoint trạng thái, và cơ chế truyền phát cập nhật theo thời gian thực. Pipeline bao gồm sáu nút (nodes) tương ứng với sáu tác tử chuyên biệt cùng một nút xử lý lỗi, được nhóm thành ba giai đoạn logic:

\begin{enumerate}
    \item \textbf{Giai đoạn Khởi tạo} (Initialization): Tác tử Giao tiếp ($\mathcal{A}_{conv}$) thu thập yêu cầu từ người dùng qua ngôn ngữ tự nhiên, trích xuất ý định dự đoán, và xác thực kết nối dữ liệu.
    \item \textbf{Giai đoạn Phân tích và Kỹ nghệ Dữ liệu} (Data Planning \& Engineering): Tác tử Phân tích ($\mathcal{A}_{eda}$) khám phá cấu trúc cơ sở dữ liệu và thực hiện phân tích dữ liệu khám phá; Tác tử Xây dựng Dữ liệu ($\mathcal{A}_{builder}$) chuyển đổi cơ sở dữ liệu quan hệ $(\mathcal{T}, \mathcal{L})$ sang đồ thị thực thể quan hệ (Mục~\ref{sec:relation-graph}); Tác tử Xây dựng Tác vụ ($\mathcal{A}_{task}$) sinh truy vấn SQL có nhúng cửa sổ thời gian để tạo bảng huấn luyện $T_{train}$ (Mục~\ref{sec:task_to_training_table}).
    \item \textbf{Giai đoạn Mô hình hóa và Thực thi} (Modeling \& Execution): Tác tử Chuyên gia GNN ($\mathcal{A}_{gnn}$) lựa chọn siêu tham số từ các nguồn tri thức ngoài và sinh kịch bản huấn luyện hiện thực hóa kiến trúc GNN thời gian (Mục~\ref{sec:encoders}); Tác tử Vận hành ($\mathcal{A}_{ops}$) thực thi huấn luyện trong môi trường cô lập và trích xuất kết quả đánh giá.
\end{enumerate}

\paragraph{Điều phối dựa trên đồ thị trạng thái.}
Trình điều phối (\texttt{PlexeOrchestrator}) xây dựng một \texttt{StateGraph} với sáu nút chức năng và một nút \texttt{error\_handler}. Các chuyển tiếp giữa các nút được kiểm soát bởi các \textit{hàm định tuyến có điều kiện} (conditional routing functions) --- ví dụ, hàm \texttt{route\_from\_conversation} kiểm tra sự tồn tại của nguồn dữ liệu và xác nhận người dùng trước khi chuyển sang giai đoạn phân tích; hàm \texttt{route\_from\_schema} xác minh rằng thư mục CSV chứa dữ liệu hợp lệ trước khi tiến đến bước xây dựng dữ liệu. Cơ chế này cho phép pipeline tự điều hướng linh hoạt: quay lại bước trước khi gặp lỗi, hoặc kết thúc sớm khi điều kiện tiên quyết không được đáp ứng, thay vì chạy tuần tự cứng nhắc. Trạng thái pipeline được lưu trữ qua \texttt{MemorySaver} checkpoint, cho phép tiếp tục phiên làm việc đã bị gián đoạn.

\paragraph{Bộ nhớ trạng thái toàn cục.}
Để đảm bảo tính nhất quán xuyên suốt quy trình, hệ thống duy trì một \textbf{bộ nhớ trạng thái toàn cục} $\mathcal{S}$ (biểu diễn qua \texttt{PipelineState} \texttt{TypedDict}), đóng vai trò là không gian lưu trữ chung giữa các tác tử. Bộ nhớ $\mathcal{S}$ bao gồm bốn nhóm thành phần:

\begin{itemize}
    \item \textbf{Lịch sử tương tác} $\mathcal{S}_{hist}$: Chuỗi các thông điệp ngôn ngữ tự nhiên giữa người dùng và hệ thống, được quản lý theo chiến lược \textit{append-only} (sử dụng \texttt{Annotated[List, operator.add]} trong LangGraph).
    \item \textbf{Tri thức nghiệp vụ} $\mathcal{S}_{schema}$: Đặc tả lược đồ cơ sở dữ liệu (\texttt{SchemaInfo}), bao gồm các ràng buộc khóa chính, khóa ngoại, nhãn thời gian, kết quả phân tích thống kê (\texttt{EDAInfo}), và bản phân loại bảng thực thể/sự kiện.
    \item \textbf{Đặc tả kỹ thuật} $\mathcal{S}_{spec}$: Ý định dự đoán của người dùng (\texttt{user\_intent}), thông tin tập dữ liệu (\texttt{DatasetInfo}) chứa dấu thời gian val/test, thông tin tác vụ (\texttt{TaskInfo}), và đường dẫn đến các tệp mã được sinh (\texttt{dataset.py}, \texttt{task.py}).
    \item \textbf{Trạng thái tiến trình} $\mathcal{S}_{exec}$: Nhật ký lỗi phân loại (\texttt{ErrorRecord}), bộ đếm thử lại có ngưỡng cấu hình riêng cho từng nút, pha hiện tại (\texttt{PipelinePhase}), và kết quả huấn luyện (\texttt{TrainingResult}).
\end{itemize}

Cơ chế bộ nhớ dùng chung theo kiến trúc \textit{blackboard}~\cite{hayes1985blackboard} loại bỏ nhu cầu giao tiếp điểm-đến-điểm (point-to-point) giữa các tác tử, từ đó giảm độ phức tạp tích hợp từ $O(n^2)$ xuống $O(n)$ khi số lượng tác tử tăng lên. Mỗi tác tử chỉ cần đọc các trường liên quan trong $\mathcal{S}$ và ghi lại kết quả đầu ra.

\paragraph{Mô hình tác tử.}
Mỗi tác tử kế thừa từ lớp trừu tượng \texttt{BaseAgent} cung cấp ba thành phần cốt lõi: (i)~một mô hình ngôn ngữ lớn (LLM) cấu hình riêng cho từng vai trò thông qua \texttt{AgentConfig}, (ii)~một tập hợp các công cụ (\textit{tools}) được đăng ký qua LangChain, và (iii)~một \textit{system prompt} nhúng tri thức chuyên gia dành riêng cho giai đoạn mà tác tử phụ trách. Kiến trúc này cho phép thay đổi mô hình nền tảng (ví dụ, GPT-4o, Claude, hoặc Gemini) cho từng tác tử riêng lẻ qua biến môi trường, mở ra khả năng tối ưu hóa chi phí-chất lượng theo vai trò. Ngoài ra, mỗi tác tử tự động tích hợp các công cụ từ các máy chủ MCP (Model Context Protocol) thông qua \texttt{MCPManager}, cho phép mở rộng bộ công cụ mà không cần sửa đổi mã nguồn tác tử.

Hình~\ref{fig:architecture} minh họa luồng điều khiển tổng thể của hệ thống.

\begin{figure}[htbp]
    \centering
    % \includegraphics[width=\linewidth]{figures/relau2ml_architecture.pdf}
    \caption{Kiến trúc tổng quan của \tool. Sáu tác tử chuyên biệt cùng một nút xử lý lỗi được tổ chức trong đồ thị trạng thái LangGraph. Mỗi tác tử đọc và ghi vào bộ nhớ trạng thái toàn cục $\mathcal{S}$; các mũi tên liền thể hiện luồng chính, các mũi tên đứt thể hiện luồng phục hồi lỗi.}
    \label{fig:architecture}
\end{figure}
% =========================================================
\section{Giai Đoạn Khởi Tạo}
% =========================================================

Giai đoạn Khởi tạo chịu trách nhiệm thu thập thông tin đầu vào và xác thực yêu cầu của người dùng trước khi hệ thống chuyển sang các bước xử lý kỹ thuật tiếp theo.

\subsection{Tác Tử Giao Tiếp ($\mathcal{A}_{conv}$, Conversational Agent)}

Tác tử Giao tiếp $\mathcal{A}_{conv}$ là giao diện tương tác duy nhất giữa người dùng và hệ thống. Nhiệm vụ của tác tử là phân tích câu lệnh ngôn ngữ tự nhiên nhằm trích xuất ý định cốt lõi cùng các tham số kết nối dữ liệu cần thiết. Khác với các hệ thống AutoML truyền thống vốn yêu cầu người dùng cung cấp cấu hình kỹ thuật tường minh (tên bảng, cột mục tiêu, loại tác vụ), $\mathcal{A}_{conv}$ cho phép người dùng diễn đạt yêu cầu bằng ngôn ngữ tự nhiên --- ví dụ: \textit{``Tôi muốn dự đoán khách hàng nào sẽ rời bỏ dịch vụ trong 30 ngày tới''} --- và tự động suy luận thông tin kỹ thuật từ đó.

\paragraph{Cơ chế hoạt động.}
Quá trình tương tác được thực hiện thông qua một vòng lặp thu thập thông tin có kiểm tra (information-gathering loop), được mô tả như sau:

\begin{enumerate}
    \item $\mathcal{A}_{conv}$ tiếp nhận yêu cầu từ người dùng và đối chiếu với tập các trường thông tin bắt buộc $\mathcal{F} = \{f_1, f_2, \ldots, f_k\}$. System prompt của tác tử định nghĩa tường minh tập $\mathcal{F}$ bao gồm: nguồn dữ liệu (đường dẫn đến thư mục CSV hoặc chuỗi kết nối cơ sở dữ liệu), mô tả mục tiêu dự đoán, và xác nhận của người dùng.
    \item Với mỗi trường $f_i \in \mathcal{F}$ còn thiếu, tác tử sinh câu hỏi làm rõ tương ứng và gửi lại người dùng. Đặc biệt, tác tử được hướng dẫn trong prompt nhận diện các trường hợp người dùng chưa rõ ràng --- ví dụ, phân biệt giữa đường dẫn tệp và mô tả dữ liệu --- để tránh hiểu sai ý định.
    \item Vòng lặp kết thúc khi toàn bộ $\mathcal{F}$ được điền đầy đủ và người dùng xác nhận kế hoạch thực hiện (bao gồm tóm tắt ý định, loại dữ liệu, và mục tiêu dự đoán).
    \item Kết quả được ghi vào $\mathcal{S}_{hist}$ (bảo toàn toàn bộ lịch sử hội thoại) và $\mathcal{S}_{spec}$ (cập nhật \texttt{user\_intent} và \texttt{data\_source}) để các giai đoạn sau truy xuất.
\end{enumerate}

\paragraph{Cơ chế chuyển tiếp.}
Sau khi $\mathcal{A}_{conv}$ hoàn tất, hàm định tuyến \texttt{route\_from\_conversation} kiểm tra ba điều kiện trước khi cho phép chuyển sang giai đoạn phân tích: (i) nguồn dữ liệu đã được xác định trong $\mathcal{S}_{spec}$, (ii) người dùng đã xác nhận kế hoạch, và (iii) không có trạng thái lỗi đang hoạt động. Nếu bất kỳ điều kiện nào không thỏa mãn, pipeline quay lại nút \texttt{conversation} để tiếp tục thu thập thông tin.
% =========================================================
\section{Giai Đoạn Phân Tích và Kỹ Nghệ Dữ Liệu}
% =========================================================

Sau khi yêu cầu được xác nhận, hệ thống phân tích cấu trúc cơ sở dữ liệu và chuẩn bị đầu vào cho mô hình hóa đồ thị. Giai đoạn này gồm ba tác tử hoạt động tuần tự, mỗi tác tử sinh một sản phẩm (artifact) mà tác tử kế tiếp sử dụng.

\subsection{Tác Tử Phân Tích ($\mathcal{A}_{eda}$)}

$\mathcal{A}_{eda}$ khám phá và trích xuất siêu dữ liệu từ cơ sở dữ liệu quan hệ. Đây là tác tử đầu tiên tương tác trực tiếp với dữ liệu thô; kết quả phân tích của nó định hình các quyết định thiết kế ở toàn bộ giai đoạn sau.

\paragraph{Quy trình phân tích.}
$\mathcal{A}_{eda}$ thực thi tuần tự ba pha qua tám công cụ chuyên dụng. \textit{Pha chuẩn bị} xác thực kết nối cơ sở dữ liệu, xuất toàn bộ bảng sang định dạng phẳng, và trích xuất siêu dữ liệu lược đồ (khóa chính, khóa ngoại, cột thời gian). \textit{Pha phân tích} thực hiện bốn phép phân tích: (i)~thống kê mô tả cho từng bảng và cột, (ii)~phát hiện vấn đề chất lượng dữ liệu theo ba mức nghiêm trọng, (iii)~xác định các cột thời gian và đề xuất ngưỡng chia tập val/test dựa trên phân vị thời gian, và (iv)~phân loại quan hệ giữa các bảng theo vai trò --- bảng dữ kiện (fact) hay bảng chiều (dimension) --- tương ứng với phân loại trong Mục~\ref{sec:relational_data}. \textit{Pha tổng hợp} gộp kết quả thành báo cáo toàn cảnh phục vụ mô hình hóa.

Kết quả được lưu thành \textbf{bản đặc tả lược đồ} $\Phi$ trong $\mathcal{S}_{schema}$, xác định vai trò của từng bảng: bảng thực thể chính, bảng ngữ cảnh, và cột thời gian tham chiếu. Phân loại fact/dimension này hướng dẫn các bước tiếp theo: bảng thực thể (dimension) được dùng làm đối tượng dự đoán, bảng sự kiện (fact) cung cấp ngữ cảnh lịch sử trong truy vấn SQL.

\subsection{Tác Tử Xây Dựng Dữ Liệu ($\mathcal{A}_{builder}$)}

$\mathcal{A}_{builder}$ tự động hóa bước chuyển đổi từ cơ sở dữ liệu quan hệ $(\mathcal{T}, \mathcal{L})$ sang đồ thị dị thể $G=(\mathcal{V},\mathcal{E},\phi,\psi)$ đã định nghĩa trong Mục~\ref{sec:relation-graph}. Sản phẩm đầu ra tuân thủ giao diện \texttt{Dataset} của RelBench~\cite{robinson2024relbench}.

\paragraph{Nguyên tắc ánh xạ.}
Quá trình xây dựng đồ thị tuân theo ba quy tắc: (i)~mỗi hàng trong bảng $T_i$ trở thành một nút $v \in \mathcal{V}$ với kiểu $\tau_v = T_i$, định danh bởi khóa chính; (ii)~mỗi quan hệ khóa ngoại giữa $T_i$ và $T_j$ trở thành cạnh có hướng, kèm cạnh ngược tương ứng với $\mathcal{L}^{-1}$ (Mục~\ref{sec:schema-graph}) để đảm bảo truyền thông điệp hai chiều; (iii)~mỗi cột được gán bộ mã hóa phù hợp với kiểu dữ liệu theo hệ thống mã hóa đa phương thức trong Mục~\ref{sec:encoders} --- các cột khóa được loại bỏ khỏi tập đặc trưng.

\paragraph{Xác thực mã sinh ra.}
Mã do $\mathcal{A}_{builder}$ sinh ra phải vượt qua cổng kiểm soát chất lượng trước khi được ghi vào tệp: kiểm tra cú pháp, xác minh sự tồn tại của lớp giao diện bắt buộc, và phát hiện thư viện thiếu khai báo import. Mã không hợp lệ bị từ chối kèm phản hồi lỗi cụ thể, theo nguyên tắc \textit{fail-fast}.

\subsection{Tác Tử Xây Dựng Tác Vụ ($\mathcal{A}_{task}$)}

$\mathcal{A}_{task}$ tự động hóa quá trình xây dựng \textbf{bảng huấn luyện} $T_{train}$ với cấu trúc $(\mathcal{K}_v, t_v, y_v)$ đã định nghĩa trong Mục~\ref{sec:task_to_training_table}. Tác tử sinh truy vấn SQL qua DuckDB để tạo bảng $\mathcal{D}_{train} = \{(e_i, y_i, t_i)\}$, trong đó $e_i$ là thực thể dự đoán, $y_i$ là nhãn tính từ dữ liệu lịch sử, và $t_i$ là mốc thời gian tham chiếu.

\paragraph{Quy trình xây dựng.}
$\mathcal{A}_{task}$ tuân theo quy trình bốn bước cố định, quy định tường minh trong system prompt: (1)~xác định loại tác vụ từ độ đo đánh giá hoặc mô tả mục tiêu; (2)~xác thực ràng buộc thời gian (khoảng cách giữa dấu thời gian val và test phải không nhỏ hơn cửa sổ dự đoán $\Delta$); (3)~phân tích cấu trúc tác vụ để chọn mẫu SQL phù hợp; (4)~sinh, kiểm thử, và đăng ký mã. Sau khi sinh mã, hệ thống xác minh loại tác vụ nhúng trong mã trùng khớp với ý định người dùng --- nếu không khớp, tệp bị xóa và tác tử thử lại, ngăn lỗi lan truyền đến giai đoạn huấn luyện.

\paragraph{Hệ thống mẫu thiết kế SQL.}
System prompt nhúng năm mẫu thiết kế SQL, mỗi mẫu phù hợp với một cấu trúc dữ liệu và ngữ nghĩa tác vụ khác nhau. Việc lựa chọn mẫu dựa trên phân tích tự động cung cấp: nguồn gốc thực thể, thống kê khoảng cách thời gian giữa các sự kiện, và tín hiệu ngữ nghĩa. Bảng~\ref{tab:sql-patterns} tóm tắt năm mẫu.

\begin{table}[htbp]
\centering
\caption{Năm mẫu thiết kế SQL cho bảng huấn luyện và điều kiện lựa chọn.}
\label{tab:sql-patterns}
\begin{tabular}{lp{4cm}p{6.5cm}}
\toprule
\textbf{Mẫu} & \textbf{Ngữ nghĩa tác vụ} & \textbf{Điều kiện lựa chọn} \\
\midrule
\textbf{A} & Churn, vắng mặt hành vi & Bảng thực thể tồn tại; kiểm tra sự tồn tại sự kiện trong cửa sổ $\Delta$ \\
\textbf{B} & Sự kiện có khoảng cách lớn & Phân vị 90 khoảng cách sự kiện $> \Delta$; cửa sổ lookback xác định từ dữ liệu \\
\textbf{C} & Dự đoán toàn thực thể & Bảng thực thể tồn tại; giá trị 0 hợp lệ (ví dụ: 0 giao dịch) \\
\textbf{D} & Thực thể có ngày tạo & Cột ngày tạo phát hiện được; lọc thực thể theo thời điểm xuất hiện \\
\textbf{Link} & Dự đoán liên kết & Tác vụ link prediction; trả về danh sách thực thể đích \\
\bottomrule
\end{tabular}
\end{table}

Quyết định giữa mẫu A và B phụ thuộc vào phân vị 90 khoảng cách sự kiện: khi giá trị này vượt $\Delta$ (ví dụ: sự kiện theo mùa, giao dịch theo quý), kiểm tra sự tồn tại trong cửa sổ $\Delta$ sẽ bỏ sót thực thể trong khoảng trống, và mẫu B với cửa sổ lookback dài hơn là lựa chọn phù hợp.

\paragraph{Phòng chống rò rỉ dữ liệu.}
Mốc thời gian $t_i$ trong bảng huấn luyện kích hoạt chiến lược lấy mẫu lân cận nhất quán theo thời gian (Mục~\ref{sec:time-consistent}): chỉ các cạnh và nút xuất hiện trước $t_i$ được sử dụng khi huấn luyện, ngăn chặn rò rỉ dữ liệu từ tương lai. Ngưỡng chia tập val/test được xác định từ phân phối thời gian do $\mathcal{A}_{eda}$ phân tích.

% =========================================================
\section{Giai Đoạn Mô Hình Hóa và Thực Thi}
% =========================================================

Giai đoạn cuối hiện thực hóa kiến trúc mạng nơ-ron đồ thị thời gian (Mục~\ref{sec:encoders}), lựa chọn siêu tham số, và điều phối huấn luyện. Giai đoạn gồm hai tác tử: $\mathcal{A}_{gnn}$ sinh kịch bản huấn luyện, $\mathcal{A}_{ops}$ thực thi và gỡ lỗi kịch bản đó.

\subsection{Tác Tử Chuyên Gia GNN ($\mathcal{A}_{gnn}$)}

$\mathcal{A}_{gnn}$ nhận định nghĩa bài toán từ $\mathcal{S}_{spec}$ cùng cấu trúc đồ thị $\mathcal{G}$, từ đó lựa chọn siêu tham số và sinh kịch bản huấn luyện hoàn chỉnh. Khác với các phương pháp tối ưu siêu tham số dựa trên huấn luyện thử nghiệm lặp (Bayesian Optimization~\cite{feurer2015autosklearn}, ASHA~\cite{liaw2018tune}), $\mathcal{A}_{gnn}$ áp dụng cơ chế \textit{tối ưu siêu tham số không cần huấn luyện} (training-free HPO) thông qua giao thức Model Context Protocol (MCP).

\paragraph{Tối ưu siêu tham số qua tổng hợp tri thức.}
$\mathcal{A}_{gnn}$ truy vấn năm nguồn tri thức dị thể qua các máy chủ MCP: (i)~quy tắc suy luận dựa trên đặc tính bài toán (số nút, số bảng, loại tác vụ), (ii)~bài báo khoa học từ Google Scholar và arXiv, (iii)~nghiên cứu có trích dẫn cao từ Semantic Scholar, và (iv)~giải pháp cuộc thi và notebook trên Kaggle. Các cấu hình thu thập được tổng hợp theo chiến lược \textit{bỏ phiếu trung vị} (ensemble median voting): với mỗi siêu tham số, hệ thống lấy giá trị trung vị từ tập đề xuất, giảm ảnh hưởng của ngoại lai từ nguồn đơn lẻ.

Quá trình truy vấn MCP được bảo vệ bởi cơ chế ngắt mạch (circuit breaker): sau ba lần thất bại liên tiếp, lời gọi đến nguồn đó bị tạm dừng. Khi tất cả nguồn không khả dụng, hệ thống quay về bộ siêu tham số mặc định.

\paragraph{Lựa chọn độ đo đánh giá.}
$\mathcal{A}_{gnn}$ ưu tiên độ đo do người dùng chỉ định; nếu không có, áp dụng bảng ánh xạ mặc định theo loại tác vụ (MAE cho hồi quy, average precision cho phân loại nhị phân, accuracy cho phân loại đa lớp, MAP cho dự đoán liên kết). Các tên gọi phổ biến (AUROC, AUC, F1) được chuẩn hóa tự động.

\paragraph{Sinh kịch bản huấn luyện từ khuôn mẫu cố định.}
Sau khi xác định siêu tham số, $\mathcal{A}_{gnn}$ sinh kịch bản huấn luyện từ một \textit{khuôn mẫu mã cố định} (code template) tuân thủ quy trình tham chiếu của RelBench~\cite{robinson2024relbench}. Kịch bản hiện thực hóa kiến trúc bốn thành phần đã mô tả trong Mục~\ref{sec:encoders}: mã hóa đặc trưng nút dị thể, mã hóa thời gian tương đối, truyền thông điệp GraphSAGE~\cite{hamilton2017inductive} trên đồ thị dị thể, và đầu dự đoán MLP. Kịch bản cũng bao gồm cấu hình hàm mất mát theo loại tác vụ, bộ nạp dữ liệu với lấy mẫu lân cận nhất quán thời gian (Thuật toán~\ref{alg:computation_graph}), vòng lặp huấn luyện với dừng sớm, và ghi kết quả đánh giá.

Thiết kế khuôn mẫu cố định --- thay vì để LLM tự sinh toàn bộ mã huấn luyện --- loại bỏ rủi ro LLM sinh mã có lỗi logic tinh vi (rò rỉ dữ liệu, sai thứ tự đánh giá, thiếu lưu mô hình). System prompt cấm tác tử viết kịch bản thủ công và bắt buộc sử dụng khuôn mẫu.

\subsection{Tác Tử Vận Hành ($\mathcal{A}_{ops}$)}
\label{sec:operation-agent}

$\mathcal{A}_{ops}$ thực thi kịch bản huấn luyện do $\mathcal{A}_{gnn}$ sinh ra. Tác tử này hoạt động theo thiết kế \textit{lai} (hybrid), lấy ý tưởng từ Self-Debugging~\cite{chen2024selfdebugging} và AutoML-Agent~\cite{trirat2025automl}: đường thực thi chính không sử dụng LLM; LLM chỉ được kích hoạt khi kịch bản thất bại.

\paragraph{Pha 1: Thực thi tất định.}
$\mathcal{A}_{ops}$ khởi tạo tiến trình con cô lập để chạy kịch bản huấn luyện, thu thập cả đầu ra chuẩn và đầu ra lỗi theo thời gian thực. Nếu kịch bản hoàn tất mà không có lỗi, tác tử chuyển trực tiếp sang trích xuất kết quả --- không tiêu tốn token LLM nào trên đường chính.

\paragraph{Pha 2: Gỡ lỗi tự động.}
Khi kịch bản thất bại, $\mathcal{A}_{ops}$ kích hoạt vòng lặp gỡ lỗi tối đa $N$ lần ($N=3$). Mỗi vòng gồm bốn bước: (1)~khôi phục toàn bộ tệp mã về bản gốc, ngăn tích lũy lỗi từ bản vá thất bại trước; (2)~phân tích traceback để xác định tệp nguồn lỗi --- có thể là tệp định nghĩa đồ thị hoặc tệp tác vụ, không nhất thiết là kịch bản huấn luyện; (3)~gửi mã lỗi, traceback, và lịch sử các lần gỡ lỗi trước cho LLM để sinh bản vá; (4)~thực thi lại kịch bản.

Nếu cả $N$ vòng thất bại, $\mathcal{A}_{ops}$ khôi phục tệp gốc trước khi trả quyền cho tầng điều phối (Mục~\ref{sec:self-recovery}). Việc khôi phục là điều kiện cần để tầng điều phối có thể sinh lại mã từ đầu mà không bị ảnh hưởng bởi bản vá hỏng.

\paragraph{Sản phẩm đầu ra.}
Khi huấn luyện hoàn tất, pipeline sinh ra năm sản phẩm: định nghĩa đồ thị, định nghĩa tác vụ, kịch bản huấn luyện, trọng số mô hình đã huấn luyện, và kết quả đánh giá. $\mathcal{A}_{ops}$ cũng sinh mã suy luận (inference code) để người dùng thực hiện dự đoán trên dữ liệu mới.

\section{Cơ chế tự phục hồi đa tầng}
\label{sec:self-recovery}

Mã nguồn do LLM sinh ra có thể chứa lỗi cú pháp, lỗi ngữ nghĩa, hoặc vi phạm ràng buộc dữ liệu. Reflexion~\cite{shinn2023reflexion} và MetaGPT~\cite{hong2023metagpt} đã cho thấy phản hồi có cấu trúc giúp LLM sửa lỗi hiệu quả hơn so với thử lại mà không cung cấp ngữ cảnh. \tool áp dụng cơ chế tự phục hồi đa tầng hoạt động ở ba cấp độ, mỗi cấp có phạm vi và chiến lược phục hồi khác nhau.

\subsection{Phân loại lỗi}

Khi một tác tử tạo ra ngoại lệ, hệ thống phân loại lỗi dựa trên kiểu ngoại lệ và từ khóa trong thông điệp. Bảng~\ref{tab:error-taxonomy} mô tả ba loại lỗi: \textit{transient} (lỗi tạm thời do mạng hoặc giới hạn API, có thể tự phục hồi khi thử lại), \textit{recoverable} (lỗi sửa được bằng phản hồi có cấu trúc, ví dụ lỗi cú pháp), và \textit{permanent} (lỗi logic nghiêm trọng không thuộc hai loại trên). Phân loại này phục vụ quyết định phục hồi có chọn lọc ở các tầng tiếp theo.

\begin{table}[htbp]
\centering
\caption{Phân loại lỗi tự động theo mức độ nghiêm trọng.}
\label{tab:error-taxonomy}
\begin{tabular}{lp{4cm}p{6cm}}
\toprule
\textbf{Loại} & \textbf{Điều kiện} & \textbf{Mô tả} \\
\midrule
Transient   & Timeout, lỗi kết nối, rate limit & Lỗi tạm thời --- có khả năng tự phục hồi khi thử lại \\
Recoverable & Lỗi cú pháp, giá trị không hợp lệ & Lỗi sửa được bằng phản hồi có cấu trúc cho LLM \\
Permanent   & (mặc định) & Lỗi logic hoặc lỗi hệ thống không tự khắc phục được \\
\bottomrule
\end{tabular}
\end{table}

\subsection{Tầng 1: Cổng xác thực và thử lại tại tác tử sinh mã}

Các tác tử sinh mã ($\mathcal{A}_{builder}$, $\mathcal{A}_{task}$) có vòng lặp tự kiểm tra: sau mỗi lần gọi LLM, tác tử kiểm tra sự tồn tại của sản phẩm mục tiêu. Nếu sản phẩm chưa được tạo, tác tử bổ sung hướng dẫn sửa lỗi vào ngữ cảnh hội thoại và thực hiện lại, tối đa $K$ lần (cấu hình riêng cho từng pha).

Trước khi ghi tệp, mã phải vượt qua cổng kiểm soát chất lượng theo nguyên tắc \textit{fail-fast}: (i)~kiểm tra cú pháp, (ii)~xác minh lớp giao diện bắt buộc (tuân thủ RelBench~\cite{robinson2024relbench}), và (iii)~phát hiện thư viện thiếu khai báo import. Mã không hợp lệ không được ghi vào tệp, ngăn lỗi lan truyền sang giai đoạn sau.

\subsection{Tầng 2: Gỡ lỗi nội bộ của Tác tử Vận hành}

Tầng này khác biệt về bản chất so với Tầng 1. Ở Tầng 1, tác tử thử lại toàn bộ quá trình sinh mã khi sản phẩm chưa tồn tại. Ở Tầng 2, mã đã hoàn chỉnh nhưng thất bại tại thời điểm thực thi (runtime error). Chiến lược phục hồi chuyển từ \textit{sinh lại} sang \textit{vá lỗi có hướng dẫn} (guided patching), được mô tả trong Mục~\ref{sec:operation-agent}.

Hai quyết định thiết kế trong vòng lặp gỡ lỗi của $\mathcal{A}_{ops}$:

\begin{itemize}
    \item \textbf{Khôi phục tệp gốc trước mỗi lần thử}: Mỗi vòng bắt đầu bằng khôi phục toàn bộ tệp về bản gốc. Không có bước này, bản vá sai ở vòng $i$ tạo lỗi mới ở vòng $i{+}1$ --- hiện tượng tích lũy lỗi khiến xác suất phục hồi giảm nhanh theo số vòng. Việc khôi phục biến các lần thử trở nên độc lập.
    \item \textbf{Quy chiếu lỗi đa tệp}: Lỗi runtime trong kịch bản huấn luyện có thể bắt nguồn từ tệp định nghĩa đồ thị hoặc tệp tác vụ. Hệ thống phân tích traceback để xác định tệp nguồn, và LLM sinh bản vá cho đúng tệp cần sửa.
\end{itemize}

Khi cả $N$ vòng thất bại, $\mathcal{A}_{ops}$ khôi phục tệp gốc trước khi trả quyền cho tầng điều phối --- điều kiện cần để Tầng 3 sinh lại mã từ đầu mà không bị ảnh hưởng bởi bản vá hỏng.

\subsection{Tầng 3: Phục hồi tại tầng điều phối}

Khi Tầng 1 hoặc Tầng 2 không phục hồi được, lỗi chuyển lên tầng điều phối. Hàm định tuyến phát hiện lỗi qua trạng thái toàn cục $\mathcal{S}$ và chuyển luồng đến nút xử lý lỗi, thực hiện hai bước: (1)~so sánh số lần thử với ngưỡng tối đa cho pha bị lỗi --- nếu vượt ngưỡng, pipeline kết thúc với trạng thái thất bại; (2)~xóa lỗi tích lũy, ghi nhật ký lỗi vào lịch sử (append-only), tăng bộ đếm thử lại, và chuyển luồng quay lại pha bị lỗi.

Đặc biệt, khi $\mathcal{A}_{ops}$ thất bại, tầng điều phối không quay lại pha thực thi mà quay về pha mô hình hóa --- $\mathcal{A}_{gnn}$ sinh lại kịch bản huấn luyện từ đầu thay vì thực thi lại kịch bản cũ.

\paragraph{Phản hồi có cấu trúc xuyên suốt.}
Tại mỗi tầng, ngữ cảnh lỗi được cung cấp cho LLM theo dạng có cấu trúc~\cite{shinn2023reflexion}: ở Tầng 1 qua thông điệp hệ thống trong lịch sử hội thoại, ở Tầng 2 qua traceback và lịch sử lỗi tích lũy, và ở Tầng 3 qua nhật ký lỗi trong $\mathcal{S}$ được tự động đưa vào ngữ cảnh khi tác tử được gọi lại.
